\section*{Chapitre \ref{ch:m3:refined-products-and-cracks} --- Produits raffinés et cracks}

\begin{solution}{exo:m3:refined-products-and-cracks:1}
Essence \qty{100.80}{\$/\barrel}, gazole \qty{117.60}{\$/\barrel}~; $(2 \times 100.80 +
117.60 - 3 \times 75)/3 = \qty{31.40}{\$/\barrel}$.
\end{solution}

\begin{solution}{exo:m3:refined-products-and-cracks:2}
7.44 barils par tonne~: $900/7.4435 = \qty{120.91}{\$/\barrel}$.
\end{solution}

\begin{solution}{exo:m3:refined-products-and-cracks:3}
La conduite estivale fait monter la demande, et l'essence de qualité estivale (\omterm{def:m3:refined-products-and-cracks:rvp}{pression de vapeur Reid} d'au plus 9.0 psi à partir du 1er
juin) exige moins de composants volatils bon marché et coûte davantage à produire. Le contrat à terme
passe à la qualité estivale au printemps, si bien que le prix monte avec la qualité.
\end{solution}

\begin{solution}{exo:m3:refined-products-and-cracks:4}
$0.47 \times 100 + 0.30 \times 110 + 0.30 \times 60 - 80 = \qty{18}{\$/\barrel}$. Les
rendements sont des volumes~: les produits sont moins denses que le brut, si bien que leur volume dépasse
celui du brut (gain de traitement).
\end{solution}

\begin{solution}{exo:m3:refined-products-and-cracks:5}
$2.62 - 2.50 - 0.08 = +0.04$ \$/gal~: la position est ouverte. Elle se clôture lorsque le cours du Golfe atteint
\qty{2.54}{\$/gal}.
\end{solution}

\begin{solution}{exo:m3:refined-products-and-cracks:6}
1\,800\,000 barils~: acheter 1\,800 CL, vendre 1\,200 RB et 600 HO.
\end{solution}

\begin{solution}{exo:m3:refined-products-and-cracks:7}
Mai 2022, \qty{62.68}{\$/\barrel}~; la moyenne de juillet 2006 à août 2026 est
\qty{21.20}{\$/\barrel}.
\end{solution}

\begin{solution}{exo:m3:refined-products-and-cracks:8}
Un raffineur est acheteur de produits et vendeur de brut. Acheter des contrats à terme sur le brut accroît son exposition au
prix du brut (il achètera de toute façon du brut, à quelque prix que ce soit) au lieu de le couvrir, et laisse ses ventes de produits
non couvertes. Il ne verrouille sa marge qu'en achetant du brut et en
\emph{vendant} des produits, dans les proportions de sa \omterm{def:m3:refined-products-and-cracks:slate}{gamme de produits raffinés}.
\end{solution}

\begin{solution}{pb:m3:refined-products-and-cracks:1}
\textbf{1.} Essence $3.2132 \times 42 = \qty{134.95}{\$/\barrel}$, gazole
\qty{178.63}{\$/\barrel}. \textbf{2.} \qty{65.61}{\$/\barrel}. \textbf{3.} Environ 3.1 fois
la moyenne de \qty{21.20}{}. \textbf{4.} La perturbation du Golfe en 2026 a retiré du marché le brut et
les produits, et l'offre de gazole était déjà tendue depuis que l'Europe a perdu
ses importations russes. \textbf{5.} Mai 2022, \qty{62.68}{\$/\barrel}.
\textbf{6.} $90\,000 \times 92 = 8.28$ millions de barils. \textbf{7.} 8\,280. \textbf{8.}
5\,520 RB et 2\,760 HO. \textbf{9.} $65.61 \times 8.28$ millions = 543.3~millions de dollars.
\textbf{10.} Le gain sur contrats à terme $(65.61 - 25) \times 8.28$ millions = 336.3~millions de dollars,
compensant la marge moindre sur les barils physiques.
\textbf{11.} Base entre son brut et le WTI Cushing (qualité et lieu). \textbf{12.}
Base entre les produits du Golfe et de New York (voir la \cref{fig:m3:refined-products-and-cracks:arb}).
\textbf{13.} Vendre les produits dans son propre ratio~: 4\,968 RB et 3\,312 HO. \textbf{14.} La
couverture reste en place pour les barils que la raffinerie ne traite plus~: elle devient une position ouverte
acheteuse de brut et vendeuse de produits, à clôturer ou à faire un roll. \textbf{15.} Marge de variation sur
les contrats à terme sur produits vendus à découvert lorsque les produits montent plus vite que le brut, sans encaissement encore issu de
la marge physique. \textbf{16.} Pour garder le potentiel de hausse si le crack s'élargit encore, pour limiter
les appels de marge, et parce que la production est incertaine. \textbf{17.} Spéculateurs, \omterm{def:m3:physical-commodity-markets:house}{maisons de négoce}
et fonds acheteurs de cracks, et les contreparties des raffineurs (producteurs de brut vendant, consommateurs de produits
achetant). \textbf{18.} Verrouiller une marge trois fois supérieure à la moyenne est attrayant, mais
la courbe est généralement en backwardation à ces moments-là, et les appels de marge sur la couverture peuvent être
importants. \textbf{19.} \emph{Résultat nommé~:} acheter 8\,280 CL, vendre 5\,520 RB et 2\,760 HO,
verrouillant \qty{65.61}{\$/\barrel} de brut, 543.3~millions de dollars pour le trimestre. \textbf{20.}
Le crack, c'est l'écart entre les produits et le brut, pas le pétrole.
\end{solution}

\begin{solution}{iq:m3:refined-products-and-cracks:1}
Trois barils de brut pour deux d'essence et un de distillat~: à peu près la \omterm{def:m3:refined-products-and-cracks:slate}{gamme de produits raffinés} d'une
raffinerie américaine, et les trois jambes sont des contrats à terme liquides. Le crack est la marge par baril de
brut.
\iqlookfor{la gamme stylisée et les jambes sur contrats à terme.}
\end{solution}

\begin{solution}{iq:m3:refined-products-and-cracks:2}
Convertir avec la densité~: le nombre de barils par tonne vaut $6.2898/\rho$ avec $\rho$ en kg/l, soit environ 7.45
pour le gazole~; il faut diviser le prix à la tonne par ce nombre.
\iqlookfor{la densité, et non un nombre magique unique.}
\end{solution}

\begin{solution}{iq:m3:refined-products-and-cracks:3}
Régresser le crack sur des indicatrices de mois calendaires, ou comparer les moyennes mensuelles par un bootstrap par blocs
sur les années~; attention au changement de qualité du contrat et au roll. Le traiter comme un
calendar spread de cracks (acheteur sur les mois d'été contre les mois d'hiver), dimensionné pour les années où il
échoue.
\iqlookfor{un test rigoureux sur des années qui se chevauchent, et un spread plutôt qu'une position directionnelle.}
\end{solution}

\begin{solution}{iq:m3:refined-products-and-cracks:4}
Un choc d'offre propre aux distillats~: arrêts de raffineries, perte d'importations (sanctions,
routes perturbées), un hiver froid, un changement de réglementation~; l'offre d'essence n'est pas touchée ou
sa demande est faible.
\iqlookfor{offre et demande propres aux produits, et non au brut.}
\end{solution}

\begin{solution}{iq:m3:refined-products-and-cracks:5}
Base de qualité et de localisation du brut, base de localisation et de qualité des produits, changements de gamme,
risque de volume (arrêts), liquidité pour les appels de marge, et le roll de la couverture d'un mois à l'autre.
\iqlookfor{base, volume et liquidité.}
\end{solution}

\begin{solution}{iq:m3:refined-products-and-cracks:6}
Données de référence par raffinerie (qualités de brut et leurs formules de prix, \omterm{def:m3:refined-products-and-cracks:slate}{gamme de produits raffinés}, coûts)~; un
entrepôt de prix avec conversions d'unités et de devises~; un moteur de marge par raffinerie calculant
la valeur des produits diminuée du coût du brut~; contrôles sur les prix manquants et les gammes périmées~; sorties par
région et un historique pour les tendances.
\iqlookfor{unités, configuration par actif, qualité des données.}
\end{solution}
